% Tamil Fake News Classification Thesis Draft
\documentclass[11pt,a4paper]{article}
\usepackage[utf8]{inputenc}
\usepackage[T1]{fontenc}
\usepackage{geometry}
\usepackage{graphicx}
\usepackage{booktabs}
\usepackage{hyperref}
\usepackage{amsmath}
\geometry{margin=1in}
\title{Automatic Detection of Fake News in Tamil News Headlines: Model, Deployment, and Evaluation}
\author{Prepared from repository artifacts}
\date{April 2026}

\begin{document}
\maketitle
\begin{abstract}
This thesis presents the design, implementation, and evaluation of an automatic fake-news detection system for Tamil news headlines. We document dataset preparation, model architecture, preprocessing, deployment as a FastAPI inference service, and recommendations for reproducible training and evaluation. The reported work uses a fine-tuned BERT-based sequence classification model (12 layers, 768 hidden size) stored in the repository under \texttt{backend/tamil_classifier/my_tamil_fake_news_model}.
\end{abstract}

\section{Introduction}
Fake news and misinformation pose significant challenges to information ecosystems. Automatic classification systems can assist fact-checkers and platforms by filtering likely false content. This work focuses on Tamil news headlines, presenting an end-to-end system: dataset, model, inference API with OCR support, and deployment artifacts.

\section{Related Work}
Prior research on fake news detection uses linguistic features, network signals, and transformer-based classifiers. Multilingual and low-resource settings frequently leverage transfer learning from multilingual BERT or language-specific BERT variants. This project follows fine-tuning of a BERT-style model for Tamil headline classification.

\section{Dataset}
The repository contains a candidate dataset: \texttt{data/Tamil-News-Headlines.csv}. The dataset should be documented for provenance, annotation methodology, and licensing. For reproducibility we recommend recording exact train/validation/test splits and random seeds. Typical dataset statistics to report include number of instances per class, average length, and class balance; placeholders are included here and should be replaced with measured values after dataset analysis.

\section{Preprocessing}
Text cleaning for inference keeps Tamil Unicode ranges and applies Indic-normalization when available (the inference service uses \texttt{indic-nlp-library}). OCR integration is provided via EasyOCR models found in \texttt{backend/tamil_classifier/ocr_models} for image-to-text extraction. For training, recommended preprocessing steps:
\begin{itemize}
  \item Normalization: Unicode normalization and Indic normalization.
  \item Tokenization: Use the provided tokenizer files in \texttt{my_tamil_fake_news_model} to match training tokenization.
  \item Truncation: Use maximum sequence length of 512 tokens, as indicated by \texttt{max_position_embeddings} in model config.
  \item Lowercasing: Follow tokenizer configuration (vocab and tokenizer files provided).
\end{itemize}

\section{Model Architecture}
The saved model follows a BERT sequence classification architecture. Extracted metadata (from \texttt{config.json}):
\begin{table}[h]
\centering
\begin{tabular}{ll}
\toprule
Property & Value \\
\midrule
Model type & BERT (BertForSequenceClassification) \\
Hidden size & 768 \\
Number of layers & 12 \\
Attention heads & 12 \\
Vocab size & 197,285 \\
Max position embeddings & 512 \\
Weights file & model.safetensors \\
Tokenizer files & tokenizer.json, vocab.txt \\
\bottomrule
\end{tabular}
\caption{Model configuration extracted from model artifact}
\end{table}

The classifier head is a standard linear layer on top of pooled BERT outputs. The repository does not include explicit training hyperparameters; see the Training section for a proposed reproducible setup.

\section{Training (Reproducible Setup)}
Training scripts are not present in the \texttt{tamil_classifier} folder. To reproduce training, use the following recommended settings as a baseline:
\begin{itemize}
  \item Optimizer: AdamW
  \item Learning rate: 2e-5 (with linear warmup and decay)
  \item Batch size: 16--32 (depending on GPU memory)
  \item Epochs: 3--5
  \item Max sequence length: 256--512
  \item Loss: CrossEntropyLoss
  \item Metrics: Accuracy, Precision, Recall, F1-score (macro and per-class)
\end{itemize}

Include logging of training and validation loss, and save checkpoints per epoch. Maintain a seed for reproducibility.

\section{Evaluation and Results}
The repository currently lacks saved evaluation metrics and confusion matrices. After reproducing training or obtaining a held-out test set, report the following:
\begin{itemize}
  \item Confusion matrix with counts and normalized rates.
  \item Precision, recall, F1 for each class (Fake/Real) and macro averages.
  \item ROC-AUC if a probability-based ranking is relevant.
  \item Error analysis: common failure modes, adversarial examples, and misclassified samples.
\end{itemize}

\section{Deployment and Inference}
The project provides a FastAPI inference service in \texttt{backend/tamil_classifier/main.py}. Key endpoints:
\begin{itemize}
  \item \texttt{/health}: health check
  \item \texttt{/predict}: JSON POST with field \texttt{text}
  \item \texttt{/predict_image_upload} and \texttt{/predict_image_url}: OCR + predict pipeline
\end{itemize}

The service loads the model via Hugging Face pipelines using the local model directory, and integrates EasyOCR for image inputs. Dockerfile and \texttt{requirements.txt} are provided for containerized deployment.

\section{Ethics and Limitations}
Discuss dataset biases, potential harms from misclassification, the need for human-in-the-loop verification for high-risk decisions, and data license/consent issues. Include statements about limitations arising from dataset provenance and the lack of documented annotation procedures.

\section{Conclusion and Future Work}
Summarize contributions: artifact-based documentation of a deployed Tamil fake-news classifier, integration of OCR and Indic normalization, and a reproducible training checklist. Future work: expand dataset size and diversity, multilingual transfer learning, model compression for mobile deployment, and improved explainability (e.g., attention-based rationales).

\section*{Appendix: Model config excerpt}
\begin{verbatim}
Model: BertForSequenceClassification
hidden_size: 768
num_hidden_layers: 12
num_attention_heads: 12
vocab_size: 197285
transformers_version: 4.57.3
\end{verbatim}

\section*{Appendix: Reproducible training command (example)}
\begin{verbatim}
python train.py \
  --model_name_or_path bert-base-uncased \
  --train_file data/Tamil-News-Headlines.csv \
  --per_device_train_batch_size 16 \
  --learning_rate 2e-5 \
  --num_train_epochs 3 \
  --max_seq_length 256 \
  --output_dir output/tamil_model
\end{verbatim}

\vspace{1cm}
\noindent Prepared from repository artifacts in April 2026.
\end{document}
